% !PRINTED 320
\weberpage{320}
%<<CONT
\weq{2}{\dfrac{(x-\alpha)(x-\beta)f'(x)}{f(x)} = a(x-\beta) + b(x-\alpha) + (x-\alpha)(x-\beta)\dfrac{f_1'(x)}{f_1(x)}.}
Die rechte Seite dieses Ausdruckes wird für $x = \alpha$ negativ, nämlich
$a(\alpha - \beta)$, und für $x = \beta$ positiv, nämlich $b(\beta - \alpha)$,
und muss also, während $x$ von $\alpha$ bis $\beta$ geht, einmal, oder
allgemeiner eine ungerade Anzahl von Malen durch Null gehen. Auf der linken
Seite von (2) haben $(x-\alpha)(x-\beta)$ und $f(x)$ ein unverändertes
Vorzeichen, und also muss $f'(x)$ eine ungerade Anzahl von Malen durch Null
gehen. Daraus folgt das Rolle'sche Theorem:

\gesperrt{XII. Sind $\alpha$ und $\beta$ zwei auf einander folgende Wurzeln
von $f(x) = 0$, so liegen zwischen $\alpha$ und $\beta$ (ohne die Grenzen
mitzurechnen) Wurzeln von $f'(x) = 0$ in ungerader Anzahl, also mindestens
eine.}

Wichtige Folgerungen ergeben sich hieraus für den Fall, dass die Gleichung
$f(x) = 0$ \gesperrt{nur reelle Wurzeln} hat. Sind diese der Grösse nach geordnet
$\alpha, \beta, \gamma, \dots$, und ist $\alpha$ eine $a$-fache, $\beta$ eine
$b$-fache, $\gamma$ eine $c$-fache etc. Wurzel, so hat $f'(x) = 0$ in
$\alpha$ eine $(a-1)$-fache, in $\beta$ eine $(b-1)$-fache, in $\gamma$ eine
$(c-1)$-fache Wurzel und zwischen $\alpha$ und $\beta$, zwischen $\beta$ und
$\gamma$ u.\,s.\,f. je mindestens eine reelle Wurzel. Die Gesammtzahl dieser
Wurzeln ist aber
\[
a + b + c + \cdots - 1 = n - 1,
\]
und da $f'(x)$ vom $(n-1)^{\mathrm{ten}}$ Grade ist, so wird diese
Minimalzahl auch nicht überschritten. Daraus folgt der Satz:

% [page alignment corrected 2026-08-07, 4x zoom strip 4/4] This fragment stopped a third of the
% way up page 320. Everything from "$f'(x) = 0$ nur reelle Wurzeln" to the end of theorem XIV is
% printed on THIS page and had been filed at the head of i0346. Moved back here.
%
% The numeral stays OUTSIDE the wrap on both, per the spec: Weber prints "XIII."/"XIV." unspaced
% and the statement spaced. XIII had its numeral inside the wrap when this fragment was written;
% corrected here at the same time.
XIII. \gesperrt{Hat $f(x) = 0$ nur reelle Wurzeln, so hat auch
$f'(x) = 0$ nur reelle Wurzeln, und von diesen sind die, die nicht
mit mehrfachen Wurzeln von $f(x)$ zusammenfallen, einfache Wurzeln, die von
den Wurzeln von $f(x)$ getrennt werden.}

Wenn wir dies Theorem wiederholt anwenden, so können wir es
folgendermaassen verallgemeinern:

XIV. \gesperrt{Hat $f(x)$ nur reelle Wurzeln, so haben auch $f'(x), f''(x),
f'''(x) \dots$ nur reelle Wurzeln. Hat eine dieser Functionen $f^{(\nu)}(x)$
eine $a$-fache Wurzel $\alpha$, und ist $a > 1$, so ist $\alpha$ eine
$(a+\nu)$-fache Wurzel von $f(x)$.}
